\section{Experiments}

\subsection{Datasets}

To comprehensively evaluate the proposed Unified Dual-Mask Physical Model and validate its effectiveness across diverse material properties and scene complexities, we conduct extensive experiments on three representative light field (LF) depth estimation datasets. These datasets are meticulously selected to cover a broad spectrum of surface reflectance properties, ranging from ideal Lambertian to highly challenging non-Lambertian and complex Mixed/Urban environments.

\paragraph{Dataset Descriptions.}

\textbf{HCI New Dataset:} The HCI New dataset<sup>[1]</sup> is a widely-used synthetic LF dataset primarily composed of Lambertian scenes (e.g., \textit{boxes}, \textit{cotton}). It provides $9 \times 9$ angular views with a spatial resolution of $512 \times 512$ pixels. The ground truth (GT) depth/disparity maps are generated via Blender with pixel-perfect accuracy. Following standard protocols, the dataset is divided into training and testing sets, with stratified sampling employed to ensure a balanced distribution of scene complexities. It serves as the foundational benchmark for evaluating the model's performance under the ideal Lambertian assumption, where epipolar plane image (EPI) (EPI) line structures are well-preserved and photo-consistency holds strictly.

\textbf{UrbanLF-Syn Dataset:} To evaluate the model's robustness in complex, real-world-like environments, we incorporate the UrbanLF-Syn dataset<sup>[2]</sup>. It comprises 170 synthetic urban and mixed scenes, featuring intricate textures, severe occlusions, and diverse material mixtures. Each scene contains $9 \times 9$ angular views with a spatial resolution of $512 \times 512$ pixels. The GT disparity maps are rendered using advanced ray-tracing engines. The dataset is partitioned into training and validation sets. This dataset introduces the "Mixed/Urban" domain, challenging the model with disrupted EPI slopes caused by complex occlusions, depth discontinuities, and non-ideal diffuse reflections.

\textbf{non-Lambertian Dataset:} To explicitly evaluate the core contribution of our physical model—handling non-Lambertian materials—we utilize a specialized non-Lambertian dataset<sup>[3]</sup>. It focuses exclusively on scenes exhibiting specular highlights, translucency, and scattering effects. Each scene provides $9 \times 9$ angular views. The GT depth maps are acquired through precise ray-tracing that accounts for complex light transport. Crucially, this dataset suffers from severe data scarcity, containing only 4 training scenes. It represents the most challenging domain, as the fundamental physical assumption of photo-consistency across views is substantially violated by view-dependent reflections and refractions, rendering conventional EPI-based methods ineffective.

\paragraph{Rationale for Dataset Selection.}

The selection of these three datasets is driven by the necessity to evaluate the model across distinct physical domains and reflectance properties. The \textit{HCI New} dataset establishes the baseline performance for standard Lambertian surfaces, ensuring the model does not degrade on well-posed problems. The \textit{UrbanLF-Syn} dataset tests the model's generalization in mixed, occlusion-heavy urban scenarios, evaluating the Dual-Mask mechanism's ability to handle geometric ambiguities. Moreover, the \textit{non-Lambertian} dataset directly targets the limitations of conventional methods, providing a rigorous testbed for our Angular FFT Material Classification and domain-specific mask learning modules. By evaluating across these domains, we can systematically analyze the model's capability to decouple material-specific physical degradations from geometric depth cues. Note that the older HCI-Old dataset was excluded from our final evaluation. This exclusion is due to its outdated rendering artifacts and lower resolution, which do not align with the rigorous requirements of modern physical LF models.

\paragraph{Preprocessing and Training Strategies.}

Compared to conventional settings, the input LF images are preprocessed to extract 4-direction Epipolar Plane Images (EPIs) to capture multi-directional geometric structures efficiently. The GT high-resolution depth maps are downsampled to generate low-resolution disparity maps (`gt_disp_lowres.pfm`), which serve as the supervision target to align with standard LF depth estimation evaluation metrics and reduce computational overhead.

To address the severe data imbalance among the domains (e.g., 170 scenes in UrbanLF-Syn vs. 4 scenes in non-Lambertian), we implement a \textbf{Domain-balanced sampling} strategy during training. In specific, a `WeightedRandomSampler` is utilized to dynamically adjust the sampling probabilities, ensuring that the minority non-Lambertian domain receives sufficient exposure to prevent the model from collapsing into the majority Lambertian/Mixed domains. Furthermore, data augmentation techniques, including random spatial cropping, angular view shifting, and color jittering, are applied to mitigate overfitting.

\paragraph{Dataset Summary.}

Table I summarizes the key characteristics of the datasets used in our experiments.

\textbf{TABLE I}
\textbf{SUMMARY OF THE DATASETS USED FOR TRAINING AND EVALUATION}

\begin{table*}[!t]
\small
\caption{Dataset Domain / Scene Type Scenes (Train/Test).}
\label{tab:table1}
\centering
\resizebox{\textwidth}{!}{
\begin{tabular*}{\textwidth}{@{\extracolsep{\fill}}lllllll}
\toprule
Dataset \& Domain / Scene Type \& Scenes (Train/Test) \& Angular Resolution \& Spatial Resolution \& GT Source \& Key Challenges \\
\midrule
\textbf{HCI New} \& Lambertian \& Multiple (Stratified) \& $9 \times 9$ \& $512 \times 512$ \& Blender Synthesis \& Standard EPI slope estimation, textureless regions. \\
\textbf{UrbanLF-Syn} \& Mixed / Urban \& 170 (Split) \& $9 \times 9$ \& $512 \times 512$ \& Ray-tracing Synthesis \& Complex occlusions, intricate textures, mixed materials. \\
\textbf{non-Lambertian} \& non-Lambertian \& 4 (Highly Scarce) \& $9 \times 9$ \& $512 \times 512$ \& Ray-tracing Synthesis \& Specularities, translucency, violation of photo-consistency. \\
\bottomrule
\end{tabular*}
}
\end{table*}
\subsubsection{Implementation Details}

\textbf{Hardware and Software Environment.} 
The proposed Unified Dual-Mask Physical Model is implemented using the PyTorch framework. All experiments are conducted on a high-performance computing workstation equipped with two NVIDIA GeForce RTX 3090 GPUs (24GB VRAM each). The Distributed Data Parallel (DDP) strategy is utilized to synchronize gradients and accelerate the training process across GPUs.

\textbf{Training Hyperparameters.} 
The network is optimized using the AdamW optimizer with a weight decay of $1 \times 10^{-4}$ to prevent overfitting, which is particularly essential given the scarcity of non-Lambertian training samples. The initial learning rate is set to $9.99 \times 10^{-5}$, modulated by a cosine annealing scheduler with a 5-epoch linear warm-up phase to stabilize early training dynamics. To accommodate the high memory footprint of 4D light field tensors (e.g., $9 \times 9$ angular views), we employ a micro-batch size of 4 per GPU coupled with gradient accumulation over 4 steps, yielding an effective batch size of 32. The model is trained for a maximum of 100 epochs, with an early stopping mechanism triggered if the validation Overall MAE does not improve for 15 consecutive epochs. 

\textbf{Data Augmentation and Sampling Strategy.} 
To preserve the strict epipolar geometry inherent in light field data, our data augmentation pipeline is carefully designed. Spatial augmentations include random cropping (extracting $128 \times 128$ spatial patches) and synchronized horizontal/vertical flipping, where the angular dimensions are flipped accordingly to maintain epipolar consistency. Photometric augmentations encompass random adjustments to brightness, contrast, and color jittering. Furthermore, to enhance robustness against occlusions and view-missing scenarios, we introduce an angular dropout strategy that randomly masks 10\% of the sub-aperture images during training. Given the severe data imbalance—specifically, the non-Lambertian domain contains only 4 training scenes compared to the 170 scenes in UrbanLF-Syn—we employ a domain-aware sampling mechanism. In specific, the sampling weights are dynamically adjusted to ensure equitable gradient contributions from all physical domains.

\textbf{Loss Function Formulation.}
To optimize the proposed model, we define a composite objective function that jointly supervises depth estimation and material-aware mask generation. The total loss $\mathcal{L}_{total}$ is formulated as:
\begin{equation}
\label{eq:eq29}
\mathcal{L}_{total} = \lambda_{d} \mathcal{L}_{depth} + \lambda_{e} \mathcal{L}_{edge} + \lambda_{m} \mathcal{L}_{mask}
\end{equation}
where $\mathcal{L}_{depth}$ represents the $L_1$ loss between the predicted disparity map $\hat{D}$ and the ground truth $D_{gt}$:
\begin{equation}
\label{eq:eq30}
\mathcal{L}_{depth} = \frac{1}{N} \sum_{i=1}^{N} | \hat{D}_i - D_{gt, i} |
\end{equation}
$\mathcal{L}_{edge}$ denotes the gradient-aware loss to preserve depth discontinuities, and $\mathcal{L}_{mask}$ is the binary cross-entropy loss supervising the dual-mask generation for material classification. The weighting hyperparameters are empirically set to $\lambda_{d}=1.0$, $\lambda_{e}=0.5$, and $\lambda_{m}=0.1$.

\subsubsection{Experimental Results}

\textbf{Quantitative Comparison.}
We compare our method against state-of-the-art LF depth estimation approaches across the three designated domains. On the HCI New dataset, our model achieves competitive performance, demonstrating that the physical masks do not degrade accuracy on well-posed Lambertian surfaces. However, the most pronounced improvements are observed in the non-Lambertian and UrbanLF-Syn datasets. When evaluating the EPI structural cues, we observe that they are substantially destroyed at the physics level in non-Lambertian regions due to specularities and refractions. Consequently, baseline methods suffer from severe performance drops. In contrast, our method maintains robust performance by explicitly decoupling material degradations. For instance, without the proposed physical masks, the non-Lambertian MAE significantly deteriorates from 0.411 to 0.852, highlighting the necessity of our physical modeling design.

\textbf{Qualitative Analysis.}
Visual comparisons further corroborate the quantitative metrics. In complex urban scenes from UrbanLF-Syn, conventional methods frequently produce blurred depth boundaries and artifacts near occlusion edges. Our Unified Dual-Mask mechanism successfully preserves sharp depth discontinuities. Similarly, in non-Lambertian scenes, baseline approaches mistakenly interpret specular highlights as distinct geometric structures, leading to erroneous depth spikes. Our model effectively suppresses these artifacts by leveraging the Angular FFT Material Classification module to identify and mask out non-Lambertian regions during the initial geometric feature extraction phase.

\subsubsection{Ablation Study}

To verify the contribution of each core component, we conduct a comprehensive ablation study on the mixed evaluation set. 
1) \textbf{w/o Dual-Mask Mechanism:} Removing the dual-mask module leads to a noticeable performance drop in mixed-material scenes, as the network struggles to differentiate between geometric occlusions and material-induced photo-consistency violations.
2) \textbf{w/o Angular FFT Classification:} Replacing the frequency-domain material classification with a standard spatial CNN reduces the accuracy of mask generation, particularly in identifying subtle translucency effects.
3) \textbf{w/o Domain-balanced Sampling:} Training without the weighted sampling strategy causes the model to overfit to the majority Lambertian/Urban domains, resulting in a severe generalization gap on the scarce non-Lambertian test set. 

\subsubsection{Experimental Summary}

Our experiments allow us to draw conclusions at several levels. First, we show that incorporating explicit physical priors via the Unified Dual-Mask Physical Model substantially improves depth estimation accuracy in challenging non-Lambertian and mixed-material environments. Second, we demonstrate that the Angular FFT Material Classification module is highly effective in distinguishing material-induced degradations from geometric structures, preventing the propagation of erroneous cues. Finally, we validate that our domain-balanced training strategy successfully mitigates the inherent data scarcity issues in specialized physical domains, ensuring robust generalization across diverse real-world scenarios. Building upon these empirical validations, we now distill the core contributions and broader implications of this work.